\documentclass[11pt]{article}
\usepackage[a4paper,margin=28mm]{geometry}
\usepackage{amsmath,amssymb,amsthm,mathtools,hyperref}
\newcommand{\ddc}{dd^c}
\newcommand{\MA}{\operatorname{MA}}
\newcommand{\PSH}{\operatorname{PSH}}
\newcommand{\Eenv}{\mathcal E_{\mathrm{env}}}
\newtheorem{theorem}{Theorem}
\newtheorem{lemma}[theorem]{Lemma}
\newtheorem{corollary}[theorem]{Corollary}
\theoremstyle{remark}\newtheorem{remark}[theorem]{Remark}
\title{Prescribed divisorial thresholds with a positive residual Hermitian metric}
\author{}\date{}
\begin{document}\maketitle
\begin{abstract}
For an effective real divisor whose subtraction leaves a positive
Hermitian metric, we identify the maximal subsolution coefficient over
the entire one-sided logarithmic class. It is the global coefficient for
the residual metric, and it is attained. For smooth positive data, the
ordinary class and the exact relative stretching class have the same
attained coefficient. A maximum-principle argument makes this coefficient
strictly smaller than the global one whenever the divisor is nonzero.
On a projective product the two coefficients have explicit values, with
ratio $(1-\lambda)^k$. These statements characterize restricted classes;
the general inverse problem for $L^1$ data is not answered here.
\end{abstract}
\noindent\textbf{Checked limited theorem.} Independent reconstruction and
final review establish the stated positive-residual results. The earlier
reconstruction rejected an input that omitted the explicit Fubini--Study
metric in the example; that counterexample and negative record are retained.
The general original S12 inverse-realization request is not answered here.

\section{Background and exact classes}
Let $X$ be connected and compact, of complex dimension $n\geq2$, let $h$
be a smooth positive Hermitian form, and let $dV$ be a smooth positive
probability. Use $\ddc=i\partial\bar\partial/(2\pi)$ and ordinary wedge
powers. All potentials have their genuine canonical quasi-psh values.
The extended measure $\MA_h^{\rm ext}(u)$ is the non-pluripolar product
with smooth-background corrections. We use the background supplied in
the original problem: for nonzero $f\geq0$ in $L^1(dV)$,
\[
 b_h(f)=\lim_{j\to\infty}c_h(\min(f,j))
\]
is the maximal coefficient of a genuine subsolution
$\MA_h^{\rm ext}(u)\geq b_h(f)f\,dV$ and is attained. This background
allows excess or infinite mass. For nonzero nonnegative bounded data,
continuous bounded existence and the unique positive normalization
scalar $c_h(f)$ follow from \cite[Theorem 0.1]{KN} and
\cite[Corollary 0.2]{Ng}. Their $\ddc$ is twice ours, so multiply
their potentials by two and absorb the smooth positive reference-volume
ratio into the density. These bounded results are not applied to
unbounded $L^1$ data.

Ordinary order $u\preceq v$ means $u\leq v+C$ for a finite constant;
$u\sim v$ means mutual order. The model envelope is
\[
 P_h[u]=\left(\sup\{v\in\PSH(X,h):v\leq0,\ v\preceq u\}\right)^*.
\]
For a genuine reference $\phi$, the class $\Eenv(X,h,\phi)$ consists
of $u\preceq\phi$ such that every $t\geq1$ admits an $h$-psh $a_t$ with
\[
 t^{-1}a_t+(1-t^{-1})\phi\leq u.
\]
Inequalities at common poles retain their canonical meaning.

Let $E=\sum_j a_j E_j$ be an effective real divisor, $a_j\geq0$,
with a smooth closed representative $\beta_E$ and a canonical logarithm
$\ell_E$, normalized by $\sup_X\ell_E=0$, such that
\[
 [E]=\beta_E+\ddc\ell_E.
\]
Locally $\ell_E=\sum_j a_j\log|s_j|^2$ plus a smooth function.
The normalization follows directly from the one-variable circle flux
of $\log|s|^2$ for our $\ddc$. Assume throughout that
\begin{equation}\label{eq:residual}
 h_E=h-\beta_E>0.
\end{equation}
In particular $\ell_E$ is genuine $h$-psh. When $E=0$, take
$\beta_E=0$ and $\ell_E=0$.

Define three restricted coefficients by taking the supremum of $c\geq0$
with $\MA_h^{\rm ext}(u)\geq cf\,dV$ over, respectively,
\[
 \begin{split}
 \mathcal A_{\leq}(E)&=\{u\in\PSH(X,h):u\preceq\ell_E\},\\
 \mathcal A_{\sim}(E)&=\{u\in\PSH(X,h):u\sim\ell_E\},\\
 \mathcal A_{\rm env}(E)&=\Eenv(X,h,\phi_E),
       \qquad \phi_E=P_h[\ell_E].
 \end{split}
\]
Denote them by $B_{\leq},B_{\sim},B_{\rm env}$; their dependence on
$h,E,f$ will be understood.

\begin{theorem}\label{th:threshold}
For nonzero nonnegative $L^1$ data,
\[
 B_{\leq}=b_{h_E}(f),
\]
and the supremum is attained by a genuine potential. The model
$\phi_E$ is genuine and $\phi_E\sim\ell_E$. If $f$ is smooth and
strictly positive, then all three coefficients are equal:
\[
 B_{\sim}=B_{\rm env}=B_{\leq}=c_{h_E}(f).
\]
They are attained by a potential of ordinary type $\ell_E$ solving the
exact finite equation. If $E\ne0$, then
\begin{equation}\label{eq:strict}
 0<c_{h_E}(f)<c_h(f)=b_h(f).
\end{equation}
\end{theorem}

\section{Canonical subtraction and actual measures}
\begin{lemma}\label{lem:subtract}
The map $w\mapsto\ell_E+w$ is a bijection from $\PSH(X,h_E)$ onto
$\mathcal A_{\leq}(E)$. It preserves the actual extended measures:
\[
 \MA_h^{\rm ext}(\ell_E+w)=\MA_{h_E}^{\rm ext}(w).
\]
Its restriction to bounded $w$ parametrizes precisely
$\mathcal A_{\sim}(E)$.
\end{lemma}
\begin{proof}
Put $S=\operatorname{Supp}E$. For $u\leq\ell_E+K$, define
$w=u-\ell_E$ first on $X\setminus S$. There $w\leq K$ and
$h_E+\ddc w=h+\ddc u\geq0$.
To extend it in a coordinate polydisk, choose a smooth quadratic $\rho$
with $\ddc\rho\geq h_E$. Choose local divisor equations $s_j$ and shrink
so that $|s|<1$, where $s$ is their product. Then
\[
 v_\epsilon=w+\rho+\epsilon\log|s|^2
\]
is psh off $S$, bounded above, and tends uniformly to $-\infty$ near
$S$. For every $M$, $\max(v_\epsilon,-M)$ is constant near $S$ and
extends psh across it. The decreasing limit as $M\to\infty$ is a
nontrivial psh extension of $v_\epsilon$. As $\epsilon\downarrow0$,
its increasing upper-regularized limit extends $w+\rho$ canonically.

The exact lower Hessian bound also survives. The extended
$R_\epsilon=\ddc v_\epsilon$ is a positive matrix-valued Radon measure.
Off $S$, $h_E+\ddc(v_\epsilon-\rho)\geq0$; on $S$ its restriction is
exactly the positive restriction of $R_\epsilon$, since the smooth
form $h_E-\ddc\rho$ gives zero mass to $S$. Splitting the coefficient
measures into $S$ and its complement proves global positivity.
Local $L^1$ convergence as $\epsilon\downarrow0$ passes this exact
bound to the canonical extension of $w$. Extensions patch because
their locally integrable distributions agree off a measure-zero set,
and their canonical representatives agree. Ordinary equivalence
gives both upper and lower bounds on $w$ off $S$, hence a bounded
canonical extension.

Conversely, a genuine $h_E$-psh $w$ has finite supremum on compact $X$,
and $u=\ell_E+w$ is genuine because
\[
 h+\ddc u=h_E+\ddc w+[E]\geq0.
\]
It satisfies $u\preceq\ell_E$. Both constructions are inverse, as
identities of locally integrable distributions and canonical functions.
On $X\setminus S$ their positive currents agree. Non-pluripolar local
products and all their smooth corrections give zero mass to $S$.
Locality therefore proves equality of the actual global measures.
No value $-\infty-(-\infty)$ is used.
\end{proof}

The maximal-subsolution characterization applied to $h_E$ now proves
$B_{\leq}=b_{h_E}(f)$. An attained residual subsolution gives an
attained one-sided subsolution through Lemma~\ref{lem:subtract}.
For general $L^1$ data this step asserts neither finite mass nor
ordinary equivalence nor stretching membership.

\section{The model and stretching membership}
\begin{lemma}\label{lem:model}
There is a finite $K\geq0$ such that
\[
 \ell_E\leq\phi_E\leq\ell_E+K.
\]
Every member of $\mathcal A_{\sim}(E)$ belongs to
$\mathcal A_{\rm env}(E)$, and the latter is contained in
$\mathcal A_{\leq}(E)$.
\end{lemma}
\begin{proof}
For an envelope competitor $v\leq0$, Lemma~\ref{lem:subtract} gives
a genuine $h_E$-psh $w=v-\ell_E$ with $w\leq-\ell_E$ off $S$.
Since $-\ell_E$ is nonnegative and integrable,
$\int_Xw^+dV\leq\int_X(-\ell_E)dV$, uniformly in $v$.
Take a finite inner/outer coordinate-ball cover and fixed smooth
quadratic corrections making $w$ psh. At a point realizing
$\sup_Xw$, submean on a fixed-radius ball, uniform bounds on the
corrections and coordinate-volume ratios give
\[
 \sup_Xw\leq C_1\int_Xw^+dV+C_2\leq K.
\]
Thus every competitor satisfies $v\leq\ell_E+K$ canonically.
The function $\ell_E$ itself is a competitor.

For completeness, finite maxima retain the precise nonclosed background
$h$. Locally convolve the common lower Hessian bound, and apply a
smooth convex regularized maximum. Its two first derivatives are
nonnegative and sum to one, and its Hessian term is positive; the
lower bound is consequently $-h$ up to an error tending uniformly to
zero. Distributional passage to the limit proves the finite-maximum
claim. Countable ball-supremum selection and increasing maxima produce
the upper-regularized envelope; after a fixed local smooth correction,
monotone psh convergence gives a genuine $h$-psh function.
It is nontrivial because of the lower competitor $\ell_E$.
The upper bound survives regularization since $\ell_E+K$ is usc.
This proves the displayed inequalities, including at $S$.

For $u\sim\ell_E\sim\phi_E$, choose $K_u\geq0$ with
$\phi_E-K_u\leq u$. The actual $h$-psh function
$a_t=\phi_E-tK_u$ satisfies
\[
 t^{-1}a_t+(1-t^{-1})\phi_E=\phi_E-K_u\leq u
\]
for every $t\geq1$, canonically at common poles. Ordinary upper order
also gives $u\preceq\phi_E$. This proves stretching membership
directly. The converse containment follows from the defining order
$u\preceq\phi_E\sim\ell_E$.
\end{proof}

For smooth positive $f$, Corollary~1 of Tosatti--Weinkove~\cite{TW}
gives a smooth $w$ with $h_E+\ddc w>0$ and
$(h_E+\ddc w)^n=c_{h_E}(f)f\,dV$. To use its normalization, take
$F=\log(f\,dV/h_E^n)$ and multiply its potential by $2\pi$.
Uniqueness of the bounded scalar identifies the source scalar.
The potential $\ell_E+w$ is of ordinary type $\ell_E$, and its actual
extended equation is finite, of mass $c_{h_E}(f)\int_Xf\,dV$.
The containments in Lemma~\ref{lem:model} and the one-sided upper bound
prove equality and attainment for all three classes.

\section{Strictness and a smooth-data obstruction}
Let $v,w$ be the smooth positive equation solutions for $h,h_E$, and put
$\sigma=h+\ddc v$, $\eta=h_E+\ddc w$. The one-sided construction
already gives $c_{h_E}(f)\leq b_h(f)=c_h(f)$.
Suppose equality holds, and set
$a=\ell_E+w-v$ on $X\setminus S$. There
$\ddc a=\eta-\sigma$ and $\eta^n=\sigma^n$. Therefore
\[
 \ddc a\wedge Q=0,\qquad
 Q=\sum_{j=0}^{n-1}\eta^j\wedge\sigma^{n-1-j}.
\]
The smooth form $Q$ defines a real second-order operator $L$ with
strictly positive principal matrix, no zeroth-order term, and $La=0$.
Since $w-v$ is smooth, $a$ tends to $-\infty$ near $S$. Its finite
maximum is attained inside some connected component $U$ of
$X\setminus S$.

We recall the elementary strong maximum argument in the form needed
here. If $z=M-a\geq0$ is $L$-harmonic and has an interior zero, but
also has a positive point, choose a ball inside $\{z>0\}$ touching
its zero set, with closure in a coordinate chart. On the annulus
between half its radius and its full radius, the barrier
\[
 b(x)=e^{-A|x-y|^2}-e^{-AR^2}
\]
satisfies $Lb>0$ for sufficiently large $A$, by ellipticity and bounded
coefficients. A small $\epsilon>0$ makes $z-\epsilon b\geq0$ on the
two boundary spheres. A negative interior minimum would have
$L(z-\epsilon b)\geq0$, contradicting $L(z-\epsilon b)<0$.
At the touching zero, the resulting inequality forces a strictly
positive inward derivative of $z$, contradicting its vanishing
gradient at a smooth interior minimum. Thus $a$ is constant on $U$.
Its boundary is nonempty and lies in $S$: otherwise $U$ would be an
open and closed subset of connected $X$, despite $S\ne\varnothing$.
Approaching that boundary contradicts $a\to-\infty$.
This proves~\eqref{eq:strict}, completing Theorem~\ref{th:threshold}.

\begin{corollary}\label{cor:smooth}
For any smooth strictly positive density $f$, no genuine threshold
subsolution $\MA_h^{\rm ext}(u)\geq b_h(f)f\,dV$ has an actual nonzero
divisorial component in $h+\ddc u$.
\end{corollary}
\begin{proof}
If $h+\ddc u\geq[E]$ for a nonzero effective divisor, replace $E$
by a sufficiently small positive multiple. Its smooth representative
then satisfies $h-\beta_E>0$. Distributionally
$w=u-\ell_E$ has $h_E+\ddc w\geq0$ and a genuine canonical
representative. The canonical identity $u=\ell_E+w$ and the finite
supremum of $w$ give $u\preceq\ell_E$. Lemma~\ref{lem:subtract} and
the maximal-subsolution characterization imply
$b_h(f)\leq b_{h_E}(f)=c_{h_E}(f)<c_h(f)=b_h(f)$, a contradiction.
\end{proof}

\section{An exact projective-product calculation}
Let $(X_0,h_0)$ be any connected compact Hermitian $m$-fold, put
$V_0=\int_{X_0}h_0^m$, and let $\kappa$ be the normalized
Fubini--Study form on $\mathbb P^k$, $k\geq2$. On their product put
\[
 N=m+k,\quad C=\binom Nm,\quad H=h_0+\kappa,\quad
 \mu=(h_0^m/V_0)\wedge\kappa^k,\quad R=\mu/dV.
\]
Take the projective hyperplane $D=\{Z_0=0\}$, also denoting its
inverse image by $D$, and set
\[
 \ell=\log\frac{|Z_0|^2}{\sum_i|Z_i|^2},\quad
 E=\lambda D,\quad\ell_E=\lambda\ell,\qquad 0<\lambda<1.
\]
Here $[D]=\kappa+\ddc\ell$ and
$H_E=h_0+(1-\lambda)\kappa>0$. The actual smooth identities
\[
 H^N=CV_0R\,dV,\qquad
 H_E^N=CV_0(1-\lambda)^kR\,dV
\]
identify the two unique bounded scalars. Theorem~\ref{th:threshold}
therefore gives
\[
 b_H(R)=CV_0,\qquad
 B_{\sim}=B_{\rm env}=B_{\leq}=CV_0(1-\lambda)^k.
\]
All restricted coefficients are attained by $u=\lambda\ell$, with
\[
 H+\ddc u=H_E+\lambda[D],\qquad
 \MA_H^{\rm ext}(u)=CV_0(1-\lambda)^kR\,dV.
\]
The arbitrary smooth reference probability is retained through $R$.

If additionally $v_-(h_0)=0$, bounded factor pullbacks and the exact
local binomial product give $v_-(H)=v_-(H_E)=0$. Indeed their masses
are $C$ and $C(1-\lambda)^k$ times the corresponding factor masses.
Both backgrounds are nonclosed, since a closed positive $h_0$ would
have every bounded mass fixed at $V_0>0$ by bounded Stokes.
This final statement is conditional on the factor hypothesis; it does
not import an unassessed example or identify a selected mass with a
global threshold.

\begin{remark}
Ordinary prescribed-type existence for nonzero $L^p$ data, $p>1$,
is already covered by Alehyane--Lu--Salouf~\cite[Theorem 5.7]{ALS}
for their model potentials. Our $\phi_E$ satisfies their model
conditions: ordinary equivalence gives
$P_h[\phi_E]=P_h[\ell_E]=\phi_E$, and
$\epsilon h-\beta_E>0$ for $\epsilon<1$ sufficiently close to one
makes $\ell_E$ an $\epsilon h$-psh competitor below $\phi_E$.
This overlaps the existence portion of Theorem~\ref{th:threshold}.
We use the older bounded equation results and canonical subtraction
in our proof, and make no literature-priority claim.
\end{remark}

\begin{remark}
The one-sided coefficient is characterized for general $L^1$ data;
ordinary and stretching attainment is asserted here only for smooth
strictly positive data. The argument supplies no nonzero divisor at
the global threshold for general $L^1$ data, no converse from an
annihilating face, and no identification of singular currents with
limits of smooth tests. No literature-priority claim is made.
\end{remark}

\begin{thebibliography}{9}
\bibitem{TW} V. Tosatti and B. Weinkove,
\emph{The complex Monge--Amp\`ere equation on compact Hermitian manifolds},
\href{https://arxiv.org/abs/0910.1390v2}{arXiv:0910.1390v2},
Corollary~1. Its smooth positive hypotheses and potential normalization
are stated at the application. The cited foundations of the named
primary results are not claimed to have been recursively certified.
\bibitem{KN} S. Ko\l odziej and N. C. Nguyen,
\emph{Weak solutions to the complex Monge--Amp\`ere equation on
Hermitian manifolds}, \href{https://arxiv.org/abs/1312.5491v2}{arXiv:1312.5491v2}, Theorem 0.1.
\bibitem{Ng} N. C. Nguyen,
\emph{The complex Monge--Amp\`ere type equation on compact Hermitian
manifolds and applications},
\href{https://arxiv.org/abs/1501.00891v1}{arXiv:1501.00891v1}, Corollary 0.2.
\bibitem{ALS} O. Alehyane, C. H. Lu and M. Salouf,
\emph{Monge--Amp\`ere equations with prescribed singularities on
compact Hermitian manifolds},
\href{https://arxiv.org/abs/2511.02339v1}{arXiv:2511.02339v1},
Definition 2.13 and Theorem 5.7. This reference records the
existence overlap, rather than an indispensable premise of our proof.
\end{thebibliography}
\end{document}
